\subsection{Эксперимент с размещением кода и параметров и сбросами между проходами}
\label{sec:c3followup}
\begin{table}[htbp]
\centering\footnotesize\setlength{\tabcolsep}{4pt}
\caption{Дополнительный опыт ESP32-C3 при одинаковом измерении бинарного решения. Каждое условие: 500 общих исходных потоков в восьми проходах со сбросами. Приведены среднее и выборочное SD восьми средних проходов; звёздочка означает SD менее $0.0001$~мкс. Средние и p95 даны в мкс; p95 --- ближайший ранг среди 4000 вызовов. Ускорение сравнивает коэффициенты и LUT при одинаковом размещении.}
\label{tab:c3_followup}
\begin{tabular}{llrrrrr}
\toprule
Код & Параметры & Коэфф. & LUT & Ускорение & Коэфф. p95 & LUT p95\\
\midrule
Flash & Flash & $5.1690\pm0.0062$ & $3.6064\pm0.0114$ & $1.433\times$ & 5.1375 & 5.9375\\
Flash & DRAM & $5.1364^{*}$ & $3.0104\pm0.0043$ & $1.706\times$ & 5.1125 & 3.0000\\
IRAM & Flash & $5.1364^{*}$ & $3.3295\pm0.0096$ & $1.543\times$ & 5.1125 & 4.6438\\
IRAM & DRAM & $5.1447\pm0.0043$ & $3.0339^{*}$ & $1.696\times$ & 5.1438 & 3.0313\\
\bottomrule
\end{tabular}
\end{table}

В дополнительном опыте использована фиксированная схема $2\times2\times2$: коэффициентное представление или LUT отсчётов, код во Flash или IRAM, параметры во Flash или DRAM. Восемь проходов с отдельными сбросами следовали порядку Williams: каждое условие встретилось в каждой позиции блока ровно один раз, как и каждая упорядоченная пара соседних разных условий. Перед 500 измеряемыми бинарными решениями выполнялся прогрев на всех 500 общих подготовленных входах. Передача по последовательному порту, инициализация, прогрев и контрольные вычисления находились за границами одиночных измерений. Общая граница измерения счётчиком циклов включала вызов ядра и бинарный порог; прерывания оставались разрешёнными. Все 4096 измерений пустого таймера равнялись одному циклу, поправка не вычиталась. Контроль счётчика циклов согласуется с заданной частотой 160~МГц.

Независимый анализ ELF подтверждает побайтное совпадение копий во Flash и IRAM для каждого семейства: 346 байт коэффициентного ядра и 206 байт ядра LUT. Поэтому сравнение размещения кода не изменяет последовательность арифметических инструкций. Два чтения счётчика циклов охватывают косвенный вызов ядра и бинарный порог; запись решения выполняется после второго чтения. Выпущенный исходный код и зафиксированные массивы параметров совпадают с артефактами загруженной сборки. Используется тот же зафиксированный GCC 8.4.0 с флагом \texttt{-Os} и отключённым LTO.

LUT оказалась быстрее при всех четырёх согласованных размещениях (табл.~\ref{tab:c3_followup}). Наименьшее среднее время, $3.0104$~мкс, получено при коде во Flash и параметрах в DRAM, против $5.1364$~мкс для коэффициентов при том же размещении. Это ускорение $1.706\times$ и сокращение времени на 41.392\%. Перенос параметров LUT из Flash в DRAM уменьшил среднее время на 16.527\% при коде во Flash и на 8.877\% при коде в IRAM. Перенос кода LUT в IRAM помог при параметрах во Flash (сокращение 7.679\%), но увеличил среднее время на 0.783\% при параметрах в DRAM. В варианте Flash/Flash LUT улучшила среднее, но имела худший p95. Перенос параметров LUT в DRAM снизил p95 до 3.0000~мкс; максимальное зарегистрированное время в этом наиболее быстром условии составило 9.0375~мкс, поэтому процентиль не является гарантией наихудшего времени исполнения.

Парный факторный анализ использует одно среднее для каждого условия и прохода, разделённого сбросами. В среднем по четырём размещениям контраст LUT минус коэффициенты равен $-1.9016~\text{мкс}$, SD восьми парных контрастов составляет $0.0040~\text{мкс}$. Перенос кода LUT из Flash в IRAM уменьшает среднее время на $0.2769~\text{мкс}$ при параметрах во Flash, но увеличивает его на $0.0236~\text{мкс}$ при параметрах в DRAM; обе направленности наблюдаются во всех восьми проходах. В приложении~\ref{app:factorial} приведены семь факторных контрастов. Эти эффекты повторных блоков описывают данную плату и протокол; они не устанавливают причину расхождения с более ранним протоколом.

Все 32\,000 бинарных решений внутри измеряемого интервала совпали с зафиксированными эталонами скомпилированных C-вычислителей. Отдельно все 32\,000 логитов из последующего неизмеряемого повторного вычисления совпали с эталонами соответствующих представлений. В каждом условии сохранились 18 ошибок классификации на 500 потоках: совпадение исполнения не тождественно точности обнаружения. Устройство также сообщило ноль несовпадений в preflight и проверках до/после блоков. Восемь проходов относятся к одной плате и фиксированной прогретой последовательности; малые SD средних не описывают вариацию между устройствами. Редкие добавки около 965 циклов сохранены без утверждения об их причине. Опыт измеряет вывод на подготовленных признаках, но не предварительную обработку, энергию, холодный старт или пиковое рабочее использование SRAM. Массивы параметров LUT в DRAM занимают 20\,554 байта, коэффициентов --- 254 байта. Результаты определяют эффективный вариант размещения в рамках общего дополнительного протокола, но не единственную причину расхождения с прежним протоколом и не новое ранжирование исторических измерений MLP/DT5.
